% engines: lualatex, pdflatex
% expect-warnings: 0
% expect-errors: 0
\documentclass{article}
\usepackage[margin=1cm]{geometry}
\usepackage{../../latex/pgf-go}
\newcommand\case[1]{\par\medskip\noindent\textbf{#1}\par\nopagebreak}
\pagestyle{empty}
\pgfgoset{view = {scale = 1}}
\begin{document}

% view/marks, view/arrows and view/lines hide each kind on its own; resume
% with marks brings the three.

\case{The goban: a mark (C3), an arrow (B2 to D4) and a line (B4 to D2)}
\begin{goban}[5, name = al]
	\stones[b]{B2}\stones[w]{D4}
	\arrow{B2}{D4}
	\line{B4}{D2}
	\marks[triangle]{C3}
\end{goban}

\case{View: all / arrows = false / lines = false / marks = false / the three false}
\pgfgoview{al}{A1:E5}
\pgfgoview[arrows = false]{al}{A1:E5}
\pgfgoview[lines = false]{al}{A1:E5}
\pgfgoview[marks = false]{al}{A1:E5}
\pgfgoview[arrows = false, lines = false, marks = false]{al}{A1:E5}

\case{resume without marks (only the stones) / with marks (the three)}
\begin{goban}[5, resume = al]\end{goban}
\begin{goban}[5, resume = {al, marks}]\end{goban}

\case{A goban with name and resume keeps what it hid: its view has no arrow, line or mark}
\begin{goban}[5, name = al2, resume = al]\end{goban}
\pgfgoview{al2}{A1:E5}

\end{document}
